\chapter{動的グリッチ表現の設計}


\section{用語および設計上の定義}

本研究では、グリッチ表現を構成する複数の視覚的特徴と、それらをライブ壁紙上でどのように変化させるかを設計対象とする。そこで、本節では、本研究で用いる主要な語句および表現要素について定義する。


\subsection{色ずれおよびRGBずれ}

\subsection{色ずれおよびRGBずれ}

本研究における「色ずれ」とは、本来同一の位置に重なって表示される色成分の位置関係が崩れ、輪郭や図形の周囲に異なる色が分離して現れる表現を指す。

特に、赤、緑、青の各色成分を異なる位置に描画することによって、本来重なって一つの色として知覚されるRGBの各成分が分離して見える状態を、本研究では「RGBずれ」と呼称する。RGBずれが生じると、画像や図形の輪郭付近に赤、緑、青の色がはみ出すように現れ、元の形状を残しながら表示のみが正常な状態から逸脱したような外観となる。

このような色成分の位置的不一致に類似した外観は、実際の表示装置や画像処理における色成分の位置ずれによっても生じ得る。また、発生原理は異なるものの、印刷物において複数の色版の位置が一致しない「版ずれ」や「見当ずれ」でも、輪郭部分に異なる色が露出する類似した視覚的特徴が現れる。このように、複数の色成分が本来の位置からずれて表示される外観は、デジタル・アナログ問わず情報伝達手段において、不整合を想起させる視覚的特徴を持つ。

RGBずれは、このような特徴を意図的に利用したグリッチ表現において広くみられる手法の一つであり、画面全体を大きく欠損させなくても、正常な表示状態からの逸脱を明確に示すことができる。そのため、本研究では帯状ノイズや欠損などと組み合わせて用いる主要なグリッチ表現の一つとして位置付ける。

本研究では、RGBずれの大きさを調整可能とし、わずかに輪郭へ色がにじむ程度の状態から、各色成分が明確に分離して見える状態まで連続的に変化させることができるよう設計する。

\subsection{欠損}

本研究における「欠損」とは、本来表示されるべき領域の一部が失われ、別の色や形状によって置き換えられている表現を指す。

本研究では、画像や画面全体を均一に変形させるのではなく、一部の領域のみを欠落または変形させることで、情報の一部が正常に描画されなかったような外観を生成する。

欠損は単独の表現として用いるだけでなく、後述する帯状ノイズ、ドット崩壊、矩形ノイズなどにも含まれる基本的な視覚要素として扱う。すなわち、本研究における欠損は特定の形状を指すものではなく、正常な表示状態の一部が局所的に失われる現象を示す概念として用いる。


\subsection{帯状ノイズ}

本研究における「帯状ノイズ」とは、画面の一部を水平方向または垂直方向の帯として区切り、その位置、色、形状などを周囲とは異なる状態、あるいはずれた状態へ変化させる表現を指す。

本研究では、特に水平方向に連続した領域をずらす表現を中心として用いる。縦方向に区切った画面の一部のみを横方向へ移動させたり、帯の内部にRGBずれや欠損を加えたりすることで、映像信号や画像データの一部が断続的に崩れたような外観を構成する。

また、帯の幅、長さ、位置、ずれ量などを変化させ、複数の帯状ノイズを重ねて用いることで、細かな線状の乱れだけでなく、画面の広い範囲が断層のように分断された状態を表現することも可能である。

このような帯状のずれや分断は、グリッチ表現において視覚的な印象を大きく左右する要素の一つであり、本研究ではグリッチらしさを形成する主要な表現として位置付ける。特に、画面の連続性を大きく崩しながらも、元の構造を部分的に残すことができるため、単純なノイズの付加とは異なる強い破損感を与えることができる。

さらに、このような断層的な表現は、デジタル画像や映像において、信号の欠落や伝送エラーなどにより情報の連続性が失われた状態を想起させるものとして利用する。


\subsection{ノイズ}

一般に「ノイズ」は、本来伝達されるべき信号や情報に対して混入する不要な成分や乱れを指す。視覚的な例としては、テレビ放送の受信状態が悪い際に画面全体へ細かな白黒の点が不規則に現れる、いわゆる「砂嵐」が代表的である。このようなノイズは、本来表示される映像とは無関係な細かな情報が画面上に大量に現れることで、正常な映像の視認を妨げる特徴を持つ。

デジタル画像や映像においても、細かな点や色の乱れ、線状の異常、局所的な画素の変化など、正常な画像情報に対して付加される不規則な視覚情報をノイズとして捉えることができる。グリッチ表現では、こうした本来であれば除去されるべき視覚的な乱れを、意図的に表現要素として利用する。

本研究における「ノイズ」とは、このような視覚的な乱れを模した表現の総称として用いる。具体的には、画面上に細かな点、粒、線、矩形などを配置し、それらの位置、大きさ、密度、色などを不均一に変化させることで構成する。

また、本研究ではノイズを画面全体へ一様に配置するのではなく、ある領域では密度を高くし、別の領域ではほとんど発生させないなど、局所的な偏りを持たせる。さらに、時間経過に応じてノイズの位置や形状を変化させることで、単なる模様として固定されるのではなく、表示状態が継続的に不安定になっているような外観を形成する。

なお、本研究では「ノイズ」を比較的広い概念として使用する。テレビの砂嵐に類似した細かな粒状の乱れだけでなく、より大きな矩形状の破片や、局所的に密集して生じる微細な崩壊なども、正常な表示に対して付加される視覚的な乱れという点ではノイズの一種として扱う。その上で、形状や大きさ、配置方法などの違いに応じて、「矩形ノイズ」や「ドット崩壊」などの名称を用いて区別する。


\subsection{矩形ノイズ}

本研究における「矩形ノイズ」とは、画面上に比較的大きな四角形や細長い矩形状の領域を配置し、局所的に表示内容が欠損、置換、あるいは変形したような印象を与える表現を指す。

細かな粒状ノイズとは異なり、矩形ノイズでは一つの領域がある程度まとまった大きさを持つ。そのため、画像や映像を構成する情報の一部分が、矩形単位で欠落または破損したような外観を生成することができる。

本研究では、矩形の大きさや縦横比を一定にせず、大型の矩形、細長い破片、小型の矩形など、異なる大きさの要素を組み合わせる。これにより、規則的な幾何学模様としてではなく、不均一なデジタルノイズとして知覚されることを目指す。


\subsection{ドット崩壊}

本研究における「ドット崩壊」とは、画面上に微細な点状または小型矩形状のノイズを多数発生させ、表示の一部が細かな粒子へ崩れているように見せる表現を指す。

ドット崩壊は画面全体に均等に配置するのではなく、主として前景側のレイヤに配置し、一部の領域では局所的に密度を高くする。これにより、細かな埃や粒子が特定の領域から発生しているような外観を形成する。

また、局所的に密集したドットとは別に、画面全体にも少量の細かなノイズを散在させる。これにより、特定領域のみが独立して崩れているように見えることを避け、画面全体に細かな表示異常が発生している状態を構成する。

本研究では、比較的大きな帯状ノイズや矩形ノイズに対して、ドット崩壊をより細かな破損を表現する要素として位置付ける。異なる大きさの破損表現を重ねることで、グリッチ表現全体の複雑さを高めることを意図している。


\subsection{フレームバッファ残像}

本研究における「フレームバッファ残像」とは、過去に描画された画面の状態を保持し、現在の描画結果と重ね合わせることによって、以前のフレームが画面上に残留しているように見せる表現を指す。

色ずれ、帯状ノイズ、矩形ノイズなどが主として画面上の空間的な破損を表現するのに対し、残像は異なる時刻の描画状態を同時に提示することで、画面の更新処理そのものに異常が発生しているような外観を構成する。

これにより、画面の一部が正常に更新されず、過去の表示状態が取り残されたような、時間的なずれや描画の停滞を表現することができる。

なお、現行の実装では、過去フレームを実際に保存して重ねる完全なフレームバッファ残像処理は用いておらず、複数の半透明レイヤや時間変化の差、色ずれなどを用いて残像に近い外観を擬似的に生成している。フレームバッファを利用した残像表現については、今後の拡張要素の一つとして検討する。


\subsection{パララックス}

パララックス（parallax）は、一般に観察位置の変化によって、対象物の見かけ上の位置が変化する「視差」を指す言葉である。例えば、観察者が左右に移動した場合、手前に存在する物体は大きく位置が変化して見える一方、遠方の物体は比較的小さく移動して見える。この位置変化の差によって、人は複数の対象物の前後関係や奥行きを知覚することができる。

この性質は、コンピュータグラフィックスやWeb表現、スマートフォンのユーザインタフェースなどにおいても利用されている。複数の画像や描画レイヤを異なる速度・移動量で動かすことによって、実際には平面的な画面上に配置された要素であっても、手前と奥に異なる距離で存在しているような擬似的な奥行きを与えることができる。このような表現は、一般にパララックス表現やパララックス効果と呼ばれる。

本研究では、このパララックスの考え方を、スマートフォン端末の傾きに応じたグリッチ表現へ応用する。画面を構成する複数のグリッチレイヤに異なる移動量を設定し、端末を傾けた際に前景側のレイヤを比較的大きく、背景側のレイヤを比較的小さく移動させることで、レイヤ間に緩やかな前後関係を与える。

本研究におけるパララックスは、正確な三次元空間を再現することを目的としたものではない。端末を傾けるという物理的な操作に対して、画面内部のグリッチ要素が異なる量で追従して移動することで、静止した平面ではなく、内部に複数の層を持つ表示面として知覚されることを意図したものである。したがって、本研究では、このような端末姿勢に応じたレイヤ間の相対移動と、それによって生じる擬似的な奥行き表現を「パララックス」と呼ぶ。

% 初期の設計では、パララックスによる奥行き感を比較的強く与え、端末を傾けることで画面内部を覗き込むような表現を目指していた。一方、制作および評価を進める中で、三次元的な奥行きが強く知覚されることで、画面そのものが破損しているという印象よりも、複数のオブジェクトが空間内に配置されている印象が強くなる可能性が指摘された。

% そこで改良後の表現では、パララックスを完全に取り除くのではなく、端末を傾けた際に各レイヤがわずかに異なる動きを示す程度に抑え、グリッチ表現の動きを補助する要素として位置付ける。また、奥行き感に対する好みには個人差があると考えられることから、パララックスの強度についても変更可能なパラメータとして扱う。

\subsection{動的表現}

本研究では、グリッチを構成する各要素が時間経過に伴って変化する性質を「動的」と呼ぶ。

すなわち、ノイズの位置や大きさ、色ずれ、帯状ノイズ、欠損などが固定された状態ではなく、時間とともに発生位置、形状、強度などを変化させる表現を「動的グリッチ表現」として扱う。

本研究における「動的」は、表示内容が時間的に変化することを広く指す概念であり、その変化の原因がユーザ入力によるものであるか否かは問わない。したがって、あらかじめ定められた映像が再生される表現や、時間経過に応じて自動的にノイズが変化する表現も、動的表現に含まれる。

一方、ユーザの操作や端末状態を入力として表示内容が変化する場合には、動的表現のうち、さらに「インタラクティブ性」を持つ表現として区別する。したがって、本研究におけるインタラクティブなグリッチ表現は、動的グリッチ表現に包含されるものとして位置付ける。

\subsection{インタラクティブ性}

「インタラクティブ（interactive）」とは、一般に、ユーザの操作や入力に対してシステムが反応し、その結果として表示内容や動作が変化する性質を指す言葉である。あらかじめ定められた内容を一方向的に提示するだけの表現とは異なり、ユーザとシステムとの間に入力と応答の関係が成立している点に特徴がある。

例えば、映像が時間経過に従って自動的に再生される場合、表示内容は変化するため動的ではあるが、ユーザの行動によってその変化が生じているわけではない。これに対して、タッチした位置に応じて表示が変化する、端末を傾けることで画面内の要素が移動する、ユーザの行動によって表示状態が変化するといった表現は、ユーザからの入力とシステム側の応答との対応を持つため、インタラクティブな表現として捉えることができる。

本研究では、このように、ユーザの操作や端末状態を入力として、画面上のグリッチ表現が変化する性質を「インタラクティブ性」と呼ぶ。

本研究では、主としてスマートフォン端末の傾きを入力として利用する。端末の傾き方向および傾き量に応じて、描画レイヤの位置、RGBずれ量、ノイズ量、帯状ノイズのずれ量、破損表現の強度などを変化させる。これにより、端末を動かすというユーザの行為と、画面上で生じる視覚的な変化とを対応させる。

したがって、本研究におけるインタラクティブ性は、単に時間経過に伴って表示内容が変化することを意味するものではない。ユーザの操作または端末の状態変化が入力となり、それに応じてグリッチ表現の状態が変化する性質を指す。



\subsection{パーソナライズ性}

本研究では、グリッチ表現を構成する色、ノイズ量、変形量、ちらつき、パララックス量などを利用者自身が選択・調整できる性質を「パーソナライズ性」と呼ぶ。

本研究におけるパーソナライズは、システムが利用者の嗜好を自動的に推定し、表示内容を自動変更することのみを意味するものではない。利用者自身が複数のプリセットから表現を選択したり、各パラメータの値を変更したりすることによって、自身の好みに適したグリッチ表現を構成できることを含めて扱う。

これにより、グリッチ表現に対する好みや視覚的刺激への許容範囲が利用者ごとに異なる場合でも、単一の表現を一律に提示するのではなく、利用者自身の選好に応じた表現へ変更できる構成を目指す。




\section{設計方針}

本研究では、スマートフォンの画面上に、単純な幾何学模様としてのノイズではなく、実際のデジタル表示やデータ処理に異常が生じた際の外観を想起させるグリッチ表現を、インタラクティブに構成することを目標とした。

前節で述べたように、グリッチ表現には、色ずれ、画像の欠損、帯状の変形、粒状ノイズ、矩形状ノイズなど、さまざまな視覚的特徴が存在する。本研究では、これらを単独のエフェクトとして提示するのではなく、複数の表現を異なる大きさ、位置、時間変化で重ね合わせることによって、画面全体に複雑な破損が発生しているような外観を構成する。

この際、本研究では、特定の一つの表現を「最も美しいグリッチ」として定めることを目的としない。グリッチ表現に対する美的評価や、ちらつき、色変化、ノイズ量などに対する許容範囲には個人差が存在すると考えられるためである。そこで、各表現要素を可能な限り独立したパラメータとして扱い、利用者自身がその種類や強さを変更できる構成とした。

例えば、強いちらつきや大きな色ずれを好む利用者はそれらの値を大きく設定できる一方、視覚的刺激を抑えたい利用者は、ちらつきを無効化したり、ノイズ量や変形量を小さくしたりすることができる。これにより、利用者の選好に応じて表現を変更できるパーソナライズ性を持たせる。

加えて、本研究では、インタラクティブなライブ壁紙としての特徴を明確にするため、端末姿勢の変化に対して視覚的な反応が生じることを重視する。グリッチ表現自体は時間経過によって継続的に変化するが、それとは別に、端末を傾けることで画面内のレイヤ位置、視差量、色ずれ、ノイズ量、破損量などが変化する構成とした。

このとき、単にすべてのグリッチを常時激しく変化させるのではなく、端末を動かした際の変化が利用者に認識できることを重視した。常時の時間変化が強すぎる場合、端末を傾けてもその影響が視覚的に埋もれてしまい、インタラクティブな表現として認識されにくくなるためである。そこで、本研究では、端末姿勢に応じて空間的な視差を生じさせる表現と、端末姿勢の変化に応じてグリッチ自体を強く発生させる表現の双方を設計した。これらについては後節で詳述する。

初期版では、まず複数レイヤによるグリッチ表現と、端末姿勢に応じた描画変化を実現することを優先した。そのため、初期版では比較的大きな矩形や単純なレイヤ移動が視覚的に目立ち、グリッチというよりも幾何学的な図形表現として知覚される部分があった。

そこで改良版では、細かな粒状ノイズや小型の矩形状ノイズ、帯状の崩壊、破片状の表現などを追加し、異なる粒度の視覚要素を重ね合わせた。これにより、大きな形状変化だけでなく、画面全体に細かな情報の乱れが存在する状態を作り、より複雑なデジタルノイズとして知覚される外観を目指した。また、利用者が表現を選択・調整しやすいよう、パラメーター調節画面においてもユーザインタフェースについて改良を行った。


\section{グリッチ表現の構成}

本研究のグリッチ表現は、複数の視覚レイヤを重ね合わせることで構成する。基礎となる構造は5層のレイヤからなり、それぞれ異なる移動量およびノイズ表現を持たせることで、画面に奥行きと複雑さを与えている。

改良の過程では、この5層に加えて、大型の矩形ノイズ、最前面に表示される細かなドット状ノイズ、画面破損を強調するオーバーレイ、および後述する「ザビ」系表現で利用する細長い破片レイヤを追加した。

描画構成は、奥側から順に、色調を形成する背景層、大型矩形ノイズ、複数の通常ノイズ層、帯状崩壊層、前景ノイズ層、最前面のドット状崩壊、細長い破片、画面破損オーバーレイによって構成される。

本システムでは、入力画像そのものを直接変形してグリッチを生成する方式を必須としていない。画像を利用する場合にはそこから色彩情報を抽出するが、実際の描画では、抽出したパレットをもとに色面およびノイズ面を生成する。そのため、元画像を使用しない場合でもグリッチ表現そのものを生成することができる。

複数のグリッチ要素を重ねる理由は、実際のデジタル表示上で生じる異常が、必ずしも単一の視覚的特徴のみとして現れるとは限らないためである。本研究では、大小のノイズ、色ずれ、帯状変形、欠損などを重ね合わせることで、単一の図形変形よりも複雑な破損感を形成することを意図した。


\section{各表現手法}

本研究では、前節で定義した複数のグリッチ表現を、単独の効果としてではなく、異なる大きさ、位置、時間変化を持つ複数のレイヤとして組み合わせて使用する。本節では、各表現の定義そのものではなく、ライブ壁紙を構成する上でそれぞれをどのような役割として配置したかについて述べる。

矩形ノイズおよびドット崩壊は、画面内に異なる粒度の破損を与える目的で使用する。比較的大きな矩形ノイズは背景側から中間層に配置し、画面構造を大きく崩す役割を持たせる。一方、ドット崩壊は主として前景側に配置し、局所的に密集した微細な破損を加える。これにより、大きな破損と細かな破損が同時に存在する状態を構成する。

帯状ノイズは、本研究におけるグリッチ表現の中心的な要素として配置する。特に画面の広い範囲を横方向に分断し、他のノイズ層とは異なるスケールで画面構造そのものをずらす役割を持たせる。また、帯内部にRGBずれや欠損を加えることで、単純な位置移動ではなく、複数の表示異常が重なっているような外観を形成する。

RGBずれは、単独の効果としてだけでなく、帯状ノイズや矩形ノイズと組み合わせて使用する。これにより、形状の破損と色成分の不一致を同時に提示し、単純な色彩変化よりも複雑な表示異常として知覚されることを意図する。

また、複数のレイヤには異なる時間変化を与える。すべてのレイヤが同一の速度やタイミングで変化すると、全体が一つのアニメーションとして知覚されやすくなるため、各レイヤの変化量や更新タイミングをずらすことで、複数の異常が独立して発生しているような状態を構成する。

さらに、改良版では、初期版と比較して細粒ノイズ、小型矩形、細長い破片などの要素を増加させた。初期版では比較的大きな図形やレイヤ移動が目立ち、幾何学的な表現として認識されやすかったためである。異なる大きさのノイズを追加することで、画面全体の情報密度を高め、より複雑なデジタル破損として知覚される外観へ変更した。

以上のように、本研究では、各グリッチ表現を個別のエフェクトとして並列に用いるのではなく、それぞれに異なる空間的・時間的役割を与え、複数のレイヤとして重ね合わせることで、全体として一つの動的グリッチ表現を構成する。


\section{動きによる表現}

本研究のグリッチ表現は、基本的に静止せず、時間経過に応じて継続的に状態を変化させる。これは、グリッチを固定された模様として提示するのではなく、表示状態そのものが不安定になっているような印象を与えるためである。

ノイズの位置や形状には疑似乱数を利用し、時間経過に応じて外観を変化させる。ただし、各フレームで完全に独立した乱数を生成するのではなく、画面上の座標、レイヤごとの値、時間、端末姿勢などを入力として疑似ランダムな値を生成する。これにより、完全に無関係な画像へ毎フレーム切り替わるのではなく、一定の構造を保ちながら不規則に変化するノイズを構成する。

また、グリッチ表現においては、高速なちらつきや急激な位置変化も表現上の特徴として利用する。静止時間が長くなると、ノイズではなく固定された模様として知覚されやすくなるため、基本的には常に何らかの時間変化を持たせている。

一方で、すべての変化を時間経過だけに依存させると、端末を傾けても表現上の差が認識されにくくなり、インタラクティブなライブ壁紙としての特徴が弱くなる。そこで本研究では、自律的な時間変化に加え、端末姿勢に応じて変化の種類や強度を変更する構成とした。


\section{端末操作とのインタラクション}

本研究では、端末の傾きに対する反応として、大きく二種類の表現方式を設計した。本研究ではこれらを「ゆら」と「ザビ」と呼ぶ。

本研究では、端末の傾きに対するグリッチ表現の反応方式として、「ゆら」と「ザビ」の2種類を設ける。これらの名称は、それぞれの動きや視覚的印象を直感的に区別するため、本研究で便宜的に付した呼称である。

「ゆら」は、日本語において物体が緩やかに揺れ動く様子を表すオノマトペ「ゆらゆら（yura-yura）」に由来する。本研究では、端末を傾けた際に複数の描画レイヤが緩やかに移動し、画面全体が揺らぐように変化する表現であることから、この名称を用いる。

「ゆら」では、端末の傾きに応じて複数の描画レイヤを異なる量だけ移動させ、パララックスを中心とした動きを生成する。基本となる5層にそれぞれ異なる移動量を設定し、前景側ほど大きく、背景側ほど小さく移動させることで、画面内に緩やかな奥行き関係を与える。

この表現は、物理的に正確な透視投影や三次元空間を再現するものではなく、複数の二次元レイヤを端末姿勢に応じて移動させることで、擬似的な奥行きと揺らぎを与えるものである。端末を傾けた際に、画面内部のグリッチ要素がそれぞれ異なる量で追従し、画面の内部に複数の層が存在するような感覚を与えることを意図している。

一方、「ザビ」は、細かなノイズや断続的なグリッチが発生する際の、ざらついた視覚的印象を表現するために本研究で用いる呼称である。「ザビザビ（zabi-zabi）」という語感から、細かなノイズが激しく発生し、画面が断続的に破損するような状態を表す名称として採用した。なお、「ザビ」は一般的なグリッチ技法を示す専門用語ではなく、本研究内で2種類の動作方式を区別するために用いる名称である。

「ザビ」では、端末を傾けた際に、グリッチ表現そのものの強度を大きく変化させる。レイヤの位置移動に加え、RGBずれ、帯状ノイズの位置ずれ、ドット崩壊、ちらつき、ノイズ量などを傾きに応じて変化させる。静止時には変化を比較的抑え、端末姿勢が一定以上変化した際にグリッチの時間変化や破損量を増加させることで、「端末を動かすことによって画面の表示状態が崩れる」ような印象を形成する。

すなわち、「ゆら」が端末姿勢の変化を主として画面内部の空間的な移動へ対応させるのに対し、「ザビ」は端末姿勢の変化を表示異常そのものの強度へ対応させる表現である。前者では緩やかなレイヤ移動と擬似的な奥行きを、後者ではノイズや破損の増加を、それぞれ傾き操作に対する主要な視覚的応答として用いる。

いずれの方式も、画面上の特定のボタンを押すことでエフェクトを直接発生させるのではなく、端末そのものを動かすことによって間接的に表示状態が変化する。これにより、明示的な操作フィードバックとしてではなく、端末を扱う行為の中で画面状態が自然に変化するインタラクションを構成する。


\section{パラメータによる表現調整}

本研究では、利用者ごとのグリッチ表現に対する選好の違いに対応するため、複数の表現要素をパラメータとして調整可能とした。

主な調整項目として、RGBずれ量、前景ノイズの大きさ、ちらつきの強度、最前面ノイズの最大サイズ、帯状ノイズの最大幅、パララックス感度、ちらつき方式、色のばらつき、アクセント色の出現頻度などを設定している。これらは基本的に独立した項目として扱われ、利用者が特定の表現要素のみを変更できるようにしている。

特に、ちらつき方式については、ちらつきを使用しない状態、常時変化する状態、端末姿勢に応じて変化する状態などを選択可能とすることで、強い刺激を好まない利用者にも対応できるようにした。同様に、RGBずれ量やパララックス感度などを小さくすることで、動きや破損感を抑えた表現へ調整することも可能である。

一方、評価実験など、すべてのパラメータを個別に操作させることが適切でない場面では、表現の特徴をまとめた簡易版のインタフェースを用いる。簡易版では、「ザビモノ」「ザビカラー」「ゆらモノ」「ゆらカラー」の4種類のプリセットから選択する構成とした。

この4種類は、端末操作に対する反応方式として「ゆら」と「ザビ」を区別し、さらに色彩表現としてモノクローム系とカラー系を組み合わせたものである。これにより、多数のパラメータを直接操作しなくても、動き方と色彩という二つの大きな特徴から表現を選択できる。

通常版では、これらのプリセットを基礎としながら個々のパラメータをさらに変更することができる。本研究では、プリセットによる大まかな選択と、詳細なパラメータ調整を組み合わせることで、操作の複雑さを抑えつつ、利用者ごとの選好に応じたグリッチ表現を構成できることを目指している。